% =====================================================================
%  From exp(-i theta/2 P) to a CZ-native quantum circuit:
%  the compilation algorithm of "improved create UCCSD circuit .py"
%
%  Companion to UCCSD_circuit_explainer.tex.  That file explains the
%  physics pipeline of arXiv:2212.08006 (operator selection, one string
%  per double, the paper's figures).  THIS file documents the general
%  compiler that turns an ordered list of Pauli-string exponentials into
%  the hardware circuit, exactly as implemented (function names below
%  match the Python code).  Compile with pdfLaTeX on Overleaf.
% =====================================================================
\documentclass[11pt]{article}
\usepackage[margin=1in]{geometry}
\usepackage{amsmath, amssymb}
\usepackage{physics}
\usepackage{booktabs}
\usepackage{xcolor}
\usepackage{hyperref}
\usepackage{listings}
\lstset{basicstyle=\ttfamily\small, columns=fullflexible}

\newcommand{\RX}{R_X}
\newcommand{\CZ}{\mathrm{CZ}}
\newcommand{\code}[1]{\texttt{#1}}

\title{From $e^{-i\frac{\theta}{2}P}$ to a CZ-native quantum circuit:\\
the compilation algorithm in \code{improved create UCCSD circuit .py}}
\author{Working notes --- every step is asserted symbolically and verified
by statevector simulation}
\date{\today}

\begin{document}
\maketitle

\begin{abstract}
Input: an ordered list of Pauli strings $P_1,\dots,P_n$ on $N$ qubits
(one per selected UCCSD double) and angles $\theta_k$.
Output: a circuit over $\{H,\ \RX(\pm\tfrac{\pi}{2}),\ \RX(\theta),\ \CZ\}$
implementing $\prod_k e^{-i\theta_k P_k/2}$, with the CZ count of the
hand-tuned circuits of arXiv:2212.08006 (H$_2$: 10, LiH: 18, F$_2$: 50).
The construction rests on a single identity --- the \emph{Pauli-frame}
formula --- plus four engineering rules: two-row hub ladders, hub
continuity, a CZ-commuting cancellation pass, and adjacent-overlap
ordering.
\end{abstract}

% =====================================================================
\section{The one identity that does all the work}
% =====================================================================
Conventions: $\RX(t)=e^{-itX/2}$, qubit $q_i$ hosts spin orbital $i$,
spin-$\alpha$ orbitals on the ``top row'' $q_0\dots q_{N/2-1}$,
spin-$\beta$ on the ``bottom row'' $q_{N/2}\dots q_{N-1}$
(so $q_i$ and $q_{i+N/2}$ are vertically adjacent on the chip).

For any Clifford $D$ and any Pauli $X_p$ on a \emph{pivot} qubit $p$,
\begin{equation}
\label{eq:frame}
  D^{\dagger}\,\RX(\theta)_p\,D
  \;=\;
  \exp\!\Big(-i\,\tfrac{\theta}{2}\, D^{\dagger} X_p D\Big).
\end{equation}
So to realise $e^{-i\frac{\theta}{2}P}$ we only need a Clifford $D$
(called the \emph{frame}) with $D^{\dagger} X_p D = \pm P$; the sign is
absorbed into $\theta$. The compiler builds $D$ from three native
conjugation rules (\code{\_conj} in the code):
\begin{center}
\begin{tabular}{lll}
\toprule
gate & rule & role \\
\midrule
$\CZ_{ab}$ & $\CZ\, X_b\, \CZ = Z_a X_b$ & \emph{hangs} a $Z$ on a neighbour\\
$H_q$      & $H Z H = X$ (and $X\!\to\!Z$, $Y\!\to\!-Y$) & \emph{opens} a qubit so more CZs can attach\\
$\RX(\pm\frac{\pi}{2})_q$ & $\RX(\tfrac{\pi}{2})^{\dagger} Z\, \RX(\tfrac{\pi}{2}) = Y$ & converts $Z \leftrightarrow Y$ (the odd-$Y$ end)\\
\bottomrule
\end{tabular}
\end{center}
Note there is \textbf{no CNOT and no $R_Z$-on-parity staircase} anywhere:
the rotation is a plain $\RX(\theta)$ on one physical qubit and all
entanglement is produced by CZs, which is exactly the native gate of the
processor in arXiv:2212.08006.

\paragraph{Reading a circuit.}
Conversely, given a circuit
$C_n \RX^{(n)} \cdots C_1 \RX^{(1)} C_0$ (Clifford layers $C_k$
alternating with rotations on pivots $p_k$), regrouping with
Eq.~\eqref{eq:frame} gives
\begin{equation}
\label{eq:regroup}
\boxed{\;
\text{circuit} = \prod_{k=n}^{1}
  \exp\!\Big(-i\tfrac{\theta_k}{2}\, D_k^{\dagger} X_{p_k} D_k\Big),
\qquad
D_k = C_{k-1}\cdots C_0,
\quad\text{provided } C_n\cdots C_0=\mathbb{1}. \;}
\end{equation}
This formula is the compiler's \emph{specification}: after all
optimisations, the code re-derives every $D_k^{\dagger}X_{p_k}D_k$
symbolically (\code{\_frame}) and asserts it equals the intended $P_k$,
and checks the residual Clifford is the identity
(\code{\_verify\_program}). Nothing is trusted; everything is checked.

% =====================================================================
\section{Compiling a single rotation (\code{\_compile\_tworow})}
% =====================================================================
Each block is emitted as
\begin{equation}
  \underbrace{B}_{\text{basis layer}}\;
  \underbrace{L}_{\text{CZ ladder}}\;
  \RX(\pm\theta)_p\;
  L^{\dagger}\; B^{\dagger},
\end{equation}
with $D = L\,B$ in Eq.~\eqref{eq:frame}. The construction has three parts.

\subsection{The two-row hub ladder $L$ (\code{\_row\_chain})}
Let $S=\{q: P_q\neq I\}$ be the support, split into the $\alpha$-row part
and the $\beta$-row part. Choose:
\begin{itemize}
\item the \textbf{pivot} $p$ = a hub qubit in the $\alpha$ part
  (median of the row support, unless \emph{hub continuity},
  Sec.~\ref{sec:multi}, overrides it);
\item the \textbf{$\beta$-hub} $h$ = the $\beta$-support qubit closest to
  $p+N/2$, i.e.\ vertically below the pivot whenever possible.
\end{itemize}
Each row is fanned into its hub by CZ/H chains running inwards from both
row ends; then the $\beta$-hub is opened with one $H$ and linked to the
pivot by a \emph{single} CZ:
\begin{align}
L \;=\;& \underbrace{\CZ_{a_1 a_2}\, H_{a_2}\, \CZ_{a_2 a_3}\cdots
        \CZ_{a_{r} p}}_{\alpha\text{-row chains into } p}\;\;
        \underbrace{\CZ_{b_1 b_2}\, H_{b_2}\cdots \CZ_{b_s h}}_{
        \beta\text{-row chains into } h}\;\;
        H_h\;\; \CZ_{p\,h}.
\end{align}
($H$ after every \emph{internal} chain node; none after terminals, none
after the pivot.) Dragging $X_p$ backwards through $L$ with the rules of
Sec.~1 yields the \emph{ladder letters}
\begin{equation}
\lambda_q=\begin{cases}
 X & q=p,\ \text{internal nodes, and the }\beta\text{-hub},\\
 Z & \text{chain terminals},
\end{cases}
\end{equation}
i.e.\ $L^{\dagger}X_pL = \pm\bigotimes_{q\in S}\lambda_q$.
On the chip every CZ is nearest-neighbour, the two rows run in
\emph{parallel}, and the rows meet through one vertical CZ
$(p,\,p{+}N/2)$ --- this is the balanced-tree layout of Figs.~13/14 and
costs depth $\mathcal{O}(\log)$ instead of a linear staircase.

\subsection{The basis layer $B$ (\code{\_basis\_for})}
The ladder letters $\lambda_q$ must be rotated into the target letters
$P_q$. For each support qubit the code brute-forces the shortest word
over $\{H,\ \RX(\pm\frac{\pi}{2})\}$ (at most two gates, eight
candidates) whose conjugation maps $P_q \to \lambda_q$ up to a $\pm$
sign. Typical outcomes:
\begin{center}
\begin{tabular}{llll}
\toprule
$\lambda_q$ & target $P_q$ & basis gates & where it appears\\
\midrule
$Z$ & $Z$ & --- & inside JW chains at terminals\\
$Z$ & $X$ & $H$ & $X$-ends of the string\\
$Z$ & $Y$ & $\RX(\frac{\pi}{2})$ & the single odd-$Y$ end\\
$X$ & $X$ & --- & internal $X$, pivot with $X$\\
$X$ & $Z$ & $H$ & JW $Z$-chain on internal nodes / hubs / pivot\\
$X$ & $Y$ & $\RX(\frac{\pi}{2})\,H$ & rare (pair mode)\\
\bottomrule
\end{tabular}
\end{center}
Any leftover $-1$ from the conjugations is returned as the block sign and
folded into the rotation angle.

\subsection{Worked example: $P_1 = Y_0X_1\,X_3Z_4X_5$ (LiH, first block)}
$\alpha$-support $\{0,1\}$, pivot $p=1$; $\beta$-support $\{3,4,5\}$,
$\beta$-hub $h=1+3=4$. Time-ordered prefix $D=BL$:
\begin{equation*}
\underbrace{\RX(\tfrac{\pi}{2})_0,\ H_3,\ H_4,\ H_5}_{B}\;\Big|\;
\underbrace{\CZ_{01}}_{\alpha}\;
\underbrace{\CZ_{34},\ \CZ_{45}}_{\beta\ \text{chains into }4}\;
\underbrace{H_4}_{\text{open hub}}\;
\underbrace{\CZ_{14}}_{\text{vertical link}}
\end{equation*}
Drag $X_1$ backwards (innermost gate first):
\begin{equation*}
X_1
\xrightarrow{\CZ_{14}} X_1Z_4
\xrightarrow{H_4} X_1X_4
\xrightarrow{\CZ_{45}} X_1X_4Z_5
\xrightarrow{\CZ_{34}} X_1Z_3X_4Z_5
\xrightarrow{\CZ_{01}} Z_0X_1Z_3X_4Z_5
\xrightarrow{H_3,H_4,H_5} Z_0X_1X_3Z_4X_5
\xrightarrow{\RX(\frac{\pi}{2})_0} Y_0X_1X_3Z_4X_5 .
\end{equation*}
Then $\RX(\theta)_1$ performs $e^{-i\frac{\theta}{2}P_1}$, and the mirror
$L^{\dagger}B^{\dagger}$ would close the block. This prefix is
\emph{identical} to the first five columns of Fig.~13 of the paper.

% =====================================================================
\section{Compiling a product of rotations}
\label{sec:multi}
% =====================================================================
Emitting (block)$\cdot$(block)$\cdots$ naively would mirror-close each
ladder and rebuild the next from scratch. Three rules recover the
paper's savings; all are general.

\subsection{Hub continuity}
The pivot/hubs of block $k{+}1$ are kept on the same qubits as block $k$
whenever they lie in the new support (\code{hub\_hint}). Because the
selected doubles share creation/annihilation indices (every F$_2$ double
contains $\hat a^{\dagger}_{11}\hat a^{\dagger}_5$; the LiH doubles share
$\hat a^{\dagger}\hat a$ pairs pairwise), consecutive strings agree on
whole row segments, and with a common hub those segments compile to
\emph{letter-for-letter identical gates} in $L_k^{\dagger}$ and $L_{k+1}$.

\subsection{The cancellation pass (\code{\_peephole}), and why CZ--CZ
commutation matters}
The optimiser repeatedly deletes adjacent inverse pairs
($H_qH_q$, $\CZ\,\CZ$ on the same pair,
$\RX(t)\RX(-t)$), where ``adjacent'' means: every gate in between
\emph{commutes} with the candidate. Sound commutation rules used:
\begin{enumerate}
\item gates on disjoint qubits commute;
\item \textbf{any two CZs commute, even sharing a qubit} (both are
  diagonal in the $Z$ basis).
\end{enumerate}
Rule 2 is the single biggest improvement over the naive pass: interface
patterns like
$\CZ_{14}\,\big[\CZ_{01}\,H_1\,\CZ_{01}\big]\,\CZ_{14}$ or
$\CZ_{34}\cdots\CZ_{34}$ straddled by gates on other qubits only cancel
once CZs may slide past each other. It alone takes LiH from 20 to 18 CZ
and F$_2$ from 58 to 50 CZ.

\paragraph{Worked interface (LiH, $P_1\!\to\!P_3$,
$P_3=Y_0Z_1X_2X_3Z_4X_5$).}
Concatenate $L_1^{\dagger}B_1^{\dagger}$ with $B_3L_3$ (same pivot 1, same
hub 4 by continuity) and cancel:
\begin{itemize}
\item $\RX(-\frac{\pi}{2})_0\,\RX(\frac{\pi}{2})_0$, $H_2H_2$, $H_3H_3$,
  $H_5H_5$ vanish ($Y_0$ and the letters on $q_2,q_3,q_5$ are unchanged);
\item $\CZ_{34}\cdots\CZ_{34}$ and $\CZ_{45}\cdots\CZ_{45}$ cancel through
  the disjoint $q_0,q_1,q_2$ gates, then $H_4 \cdots H_4$ cancels, then
  the whole shared $\beta$ subtree is gone;
\item what survives is exactly the paper's interface
\begin{equation*}
  C_1 \;=\; \CZ_{14}\;\; \CZ_{01}\;\; H_1\;\; \CZ_{01}\;\; \CZ_{12}\;\;
  \CZ_{14},
\end{equation*}
  read: re-hang $Z_4$, toggle the pivot letter $X_1\!\to\!Z_1$ by the
  $\CZ_{01} H_1 \CZ_{01}$ sandwich \emph{without releasing} $Y_0$, hang
  the new $X_2$ ($H_2$ was pre-placed in the very first basis layer), and
  re-link the hub. The two $\CZ_{14}$ survive because $H_1$ stands
  between them --- precisely what Fig.~13 shows.
\end{itemize}

\subsection{Ordering (\code{\_auto\_order})}
The interfaces are cheap only between strings that overlap. The compiler
brute-forces the order of the (few) doubles to maximise
$\sum_k |\{q: P^{(k)}_q = P^{(k+1)}_q \neq I\}|$. For LiH, given
Eq.~(23)'s order $(P_1,P_2,P_3)$ with overlaps $(2,4)$, it finds the
figure's order $(P_1,P_3,P_2)$ with overlaps $(4,4)$. For the nested
F$_2$ strings the given order is already optimal.

\subsection{Optional exact pair (\code{pair=True})}
For one double, the two-string ``qubit excitation''
\begin{equation}
e^{-i\frac{\theta}{2}\,Y_aX_bX_cX_e\cdots}\;
e^{+i\frac{\theta}{2}\,X_aY_bX_cX_e\cdots}
=\exp\Big(\!-i\tfrac{\theta}{2}\big(Y_aX_b-X_aY_b\big)X_cX_e\cdots\Big)
\end{equation}
(the factors commute) realises a determinant-pair Givens rotation with no
spurious phases on \emph{any} input state. Fig.~12 (H$_2$) is exactly
this pair; the interface between the two strings only swaps
$Y_aX_b \leftrightarrow X_aY_b$, so the shared remainder cancels and the
pair costs 10 CZ instead of $2\times 6$. The default single string (6 CZ)
is variationally equivalent when acting on the reference determinant.

% =====================================================================
\section{The full algorithm and its guarantees}
% =====================================================================
\begin{enumerate}
\item map each double $(k,l,i,j)$ to its odd-$Y$ JW representative
  (\code{jw\_string\_for\_double}); optionally add the partner string
  (\code{pair});
\item order the strings by adjacent overlap (\code{\_auto\_order});
\item for each string: build $B$, $L$ with the two-row hub layout and hub
  continuity (\code{\_compile\_tworow}), emit
  $B\,L\,\RX(\pm\theta)\,L^{\dagger}B^{\dagger}$;
\item run the CZ-commuting cancellation pass (\code{\_peephole});
\item verify (\code{\_verify\_program}): every rotation's frame
  $D_k^{\dagger}X_{p_k}D_k$ equals its intended $\pm P_k$, and the
  Clifford product of the whole program is the identity
  (Eq.~\eqref{eq:regroup});
\item emit the Qiskit circuit.
\end{enumerate}

Results (statevector-verified against
$\prod_k e^{-i s_k\theta_k P_k/2}$ to $\sim\!10^{-15}$):
\begin{center}
\begin{tabular}{lcccc}
\toprule
 & this compiler & paper (Figs.~12--14) &
 \begin{tabular}{@{}c@{}}8 staircases\\per selected double\end{tabular} &
 \begin{tabular}{@{}c@{}}full textbook UCCSD\\(paper's count)\end{tabular}\\
\midrule
H$_2$ (pair / single) & 10 / 6 & 10 & 48  & 56\\
LiH                   & 18     & 18 & 208 & 280\\
F$_2$                 & 50     & 50 & 560 & 2920\\
\bottomrule
\end{tabular}
\end{center}

\paragraph{Relation to the standard construction.}
The CNOT staircase of arXiv:2005.14475 is the special case of
Eq.~\eqref{eq:frame} with a linear ladder, $R_Z$ pivot and
CNOT $= (I\otimes H)\,\CZ\,(I\otimes H)$; it spends $2(m{-}1)$
entanglers per weight-$m$ string and uncomputes everything between
rotations. The improvements here --- one string per double, CZ-native
frames, row-parallel hub trees, shared-segment reuse --- are what close
the gap from 48 to 6--10 CZs per double while keeping the circuit
mathematically identical to the advertised product of exponentials.

\end{document}
